
\documentclass{article}
\usepackage{colortbl}
\usepackage{makecell}
\usepackage{multirow}
\usepackage{supertabular}

\begin{document}

\newcounter{utterance}

\centering \large Interaction Transcript for game `wordle\_withclue', experiment `high\_frequency\_words\_clue\_no\_critic', episode 4 with reflexion{-}Qwen3.5{-}27B.
\vspace{24pt}

{ \footnotesize  \setcounter{utterance}{1}
\setlength{\tabcolsep}{0pt}
\begin{supertabular}{c@{$\;$}|p{.15\linewidth}@{}p{.15\linewidth}p{.15\linewidth}p{.15\linewidth}p{.15\linewidth}p{.15\linewidth}}
    \# & \multicolumn{2}{c}{Player} && \multicolumn{2}{c}{Game Master} \\
    \hline

    \theutterance \stepcounter{utterance}  
    & & & \multicolumn{4}{p{0.6\linewidth}}{
        \cellcolor[rgb]{0.9,0.9,0.9}{
            \makecell[{{p{\linewidth}}}]{
                \texttt{\tiny{[P1$\langle$GM]}}
                \texttt{You have attempted this same task before and did not succeed. These are the plans you wrote after those attempts.} \\
\texttt{Trial \#1: New plan: The previous attempt was aborted because the response violated the strict formatting rule "Do not generate any other text" by including a long, rambling internal monologue and failing to provide a valid "guess:" line with a five{-}letter word. The explanation became a stream of consciousness rather than a concise summary. Next time, I will strictly adhere to the required format: a single line starting with "explanation:" containing a brief, logical deduction based on the clue and feedback, followed immediately by a single line starting with "guess:" containing exactly one valid five{-}letter English word. I will ensure the explanation does not contain conversational filler, self{-}correction, or incomplete thoughts, and I will verify the word length before submitting.} \\
\\ 
\texttt{Trial \#2: New plan: The previous attempt failed because the response violated the strict formatting rule "Do not generate any other text" by including a long, rambling internal monologue and failing to provide a valid "guess:" line with a five{-}letter word. The explanation became a stream of consciousness rather than a concise summary. Next time, I will strictly adhere to the required format: a single line starting with "explanation:" containing a brief, logical deduction based on the clue and feedback, followed immediately by a single line starting with "guess:" containing exactly one valid five{-}letter English word. I will ensure the explanation does not contain conversational filler, self{-}correction, or incomplete thoughts, and I will verify the word length before submitting. Specifically, for the clue "1985 hit by heart," I will identify the song "What About Love" or the band name "Heart" and select a valid 5{-}letter candidate like "ABOUT" or "HEART" (if allowed as a word) or a related term, ensuring the output is clean and strictly follows the `explanation:` and `guess:` tags without extra text.} \\
\\ 
\\ 
\texttt{You are a language wizard who likes to guess words by using the given rules.} \\
\\ 
\texttt{Welcome to Wordle! You have six attempts to guess the target word, a valid English word of five lowercase letters (a{-}z). Please use the tags "explanation:" and "guess:" to provide a concise explanation for each guess.} \\
\\ 
\texttt{To help you make an informed guess, you will receive a clue for the word, such as} \\
\texttt{clue: snowy white.} \\
\\ 
\texttt{Here is an example guess based on the clue:} \\
\texttt{explanation: In the fairy tail Snow White, the girl is killed because she eats a poisoned apple. And the word apple has 5 letters.} \\
\texttt{guess: apple} \\
\\ 
\texttt{After each guess, your answer will be validated, and you will receive feedback indicating which letters are correct (green), which letters are correct but in the wrong position (yellow), and which letters are incorrect (red). This feedback can be useful in determining which letters to include or exclude in your next guess.} \\
\\ 
\texttt{For example, the feedback for "apple" might be:} \\
\texttt{guess\_feedback: a<yellow> p<yellow> p<green> l<yellow> e<red>} \\
\\ 
\texttt{The explanation should contain details about how both the guess\_feedback and clue work together in tandem to formulate a new guess.} \\
\\ 
\texttt{Do not generate any other text.} \\
\\ 
\texttt{Let's begin with your first guess.} \\
\\ 
\texttt{clue: 1985 hit by heart} \\
            }
        }
    }
     \\ \\

    \theutterance \stepcounter{utterance}  
    & \multicolumn{4}{p{0.6\linewidth}}{
        \cellcolor[rgb]{0.561,0.737,0.561}{
            \makecell[{{p{\linewidth}}}]{
                \texttt{\tiny{[P1$\rangle$GM]}}
                \texttt{explanation: The clue refers to the 1985 hit song "What About Love" by the band Heart, and the word "ABOUT" is a five{-}letter word contained within the song title.} \\
\texttt{guess: about} \\
            }
        }
    }
    & & \\ \\

    \theutterance \stepcounter{utterance}  
    & & & \multicolumn{2}{p{0.3\linewidth}}{
        \cellcolor[rgb]{0.9,0.9,0.9}{
            \makecell[{{p{\linewidth}}}]{
                \texttt{\tiny{[GM$|$GM]}}
                \texttt{attempts = 1} \\
\texttt{target\_word = never} \\
\texttt{target\_word\_clue = 1985 hit by heart} \\
\texttt{guess = about} \\
\texttt{guess\_feedback = a<red> b<red> o<red> u<red> t<red>} \\
            }
        }
    }
    & & \\ \\

    \theutterance \stepcounter{utterance}  
    & & & \multicolumn{4}{p{0.6\linewidth}}{
        \cellcolor[rgb]{0.9,0.9,0.9}{
            \makecell[{{p{\linewidth}}}]{
                \texttt{\tiny{[P1$\langle$GM]}}
                \texttt{guess\_feedback: a<red> b<red> o<red> u<red> t<red>} \\
\\ 
\texttt{Provide your response only in this format.} \\
\texttt{explanation: details} \\
\texttt{guess: word} \\
            }
        }
    }
     \\ \\

    \theutterance \stepcounter{utterance}  
    & \multicolumn{4}{p{0.6\linewidth}}{
        \cellcolor[rgb]{0.561,0.737,0.561}{
            \makecell[{{p{\linewidth}}}]{
                \texttt{\tiny{[P1$\rangle$GM]}}
                \texttt{explanation: Since "ABOUT" was entirely incorrect, I will try the band name "HEART" itself, which is a valid five{-}letter word and directly referenced in the clue.} \\
\texttt{guess: heart} \\
            }
        }
    }
    & & \\ \\

    \theutterance \stepcounter{utterance}  
    & & & \multicolumn{2}{p{0.3\linewidth}}{
        \cellcolor[rgb]{0.9,0.9,0.9}{
            \makecell[{{p{\linewidth}}}]{
                \texttt{\tiny{[GM$|$GM]}}
                \texttt{attempts = 2} \\
\texttt{target\_word = never} \\
\texttt{target\_word\_clue = 1985 hit by heart} \\
\texttt{guess = heart} \\
\texttt{guess\_feedback = h<red> e<green> a<red> r<yellow> t<red>} \\
            }
        }
    }
    & & \\ \\

    \theutterance \stepcounter{utterance}  
    & & & \multicolumn{4}{p{0.6\linewidth}}{
        \cellcolor[rgb]{0.9,0.9,0.9}{
            \makecell[{{p{\linewidth}}}]{
                \texttt{\tiny{[P1$\langle$GM]}}
                \texttt{guess\_feedback: h<red> e<green> a<red> r<yellow> t<red>} \\
\\ 
\texttt{Provide your response only in this format.} \\
\texttt{explanation: details} \\
\texttt{guess: word} \\
            }
        }
    }
     \\ \\

    \theutterance \stepcounter{utterance}  
    & \multicolumn{4}{p{0.6\linewidth}}{
        \cellcolor[rgb]{0.561,0.737,0.561}{
            \makecell[{{p{\linewidth}}}]{
                \texttt{\tiny{[P1$\rangle$GM]}}
                \texttt{explanation: The feedback shows 'e' is correct in the second position and 'r' is present but not in the fourth position. The song title "What About Love" contains the word "LOVE", which fits the pattern \_E\_R\_ if we consider "HEART" was wrong, but "LOVE" is 4 letters. Re{-}evaluating the song "What About Love", the word "LOVE" is key. However, looking for a 5{-}letter word with E in position 2 and R elsewhere, "WHERE" fits the pattern \_E\_R\_ (W{-}H{-}E{-}R{-}E) but H was red. Let's try "HEARD" (H was red). Let's try "EARTH" (E is 2nd, R is 4th {-} but R was yellow in HEART meaning it's not 4th). Wait, in HEART, R was yellow at pos 4, so R is not at 4. E is green at 2. So pattern is \_E\_R\_ or \_E\_\_R. The song is "What About Love". Maybe the word is "LOVED"? No. "HEART" is the band. "1985 hit". "What About Love". "Magic Man" (1976). "Barracuda" (1977). "Crazy on You" (1976). "What About Love" (1985). "These Dreams" (1986). "Alone" (1987). The word might be "ALONE"? A was red. "DREAM"? D, R, E, A, M. A was red. "MAGIC"? M, A, G, I, C. A was red. "BARRA"? No. "CRAZY"? C, R, A, Z, Y. A was red. "ABOUT" was all red. "HEART" gave E green (2), R yellow (4). So R is in the word but not at 4. E is at 2. The word "WHERE" has E at 2, R at 4 (bad). "THERE" has E at 2, R at 4 (bad). "NEVER" has E at 2, R at 5. N, E, V, E, R. E is at 2 and 4. R is at 5. In HEART, R was at 4 (yellow), so R is not at 4. In NEVER, R is at 5. This works. Is "NEVER" related? "What About Love" lyrics? "I never thought I'd see the day". Or maybe the word is "LOVER"? L, O, V, E, R. O was red in ABOUT. So no. "LOVED} \\
            }
        }
    }
    & & \\ \\

    \theutterance \stepcounter{utterance}  
    & & & \multicolumn{2}{p{0.3\linewidth}}{
        \cellcolor[rgb]{0.9,0.9,0.9}{
            \makecell[{{p{\linewidth}}}]{
                \texttt{\tiny{[GM$|$GM]}}
                \texttt{Error: The response should contain the 'guess:' keyword exactly once.} \\
            }
        }
    }
    & & \\ \\

    \theutterance \stepcounter{utterance}  
    & & & \multicolumn{2}{p{0.3\linewidth}}{
        \cellcolor[rgb]{0.9,0.9,0.9}{
            \makecell[{{p{\linewidth}}}]{
                \texttt{\tiny{[GM$|$GM]}}
                \texttt{game\_result = ABORT} \\
            }
        }
    }
    & & \\ \\

    \theutterance \stepcounter{utterance}  
    & & \multicolumn{4}{p{0.6\linewidth}}{
        \cellcolor[rgb]{0.95,0.95,0.95}{
            \makecell[{{p{\linewidth}}}]{
                \texttt{\tiny{[P1$\leftrightarrow$P]}}
                \texttt{Post{-}episode Reflexion (not sent to GM):} \\
\texttt{New plan: The previous attempt failed because the response violated the strict formatting rule "Do not generate any other text" by including a long, rambling internal monologue and failing to provide a valid "guess:" line with a five{-}letter word. The explanation became a stream of consciousness rather than a concise summary. Next time, I will strictly adhere to the required format: a single line starting with "explanation:" containing a brief, logical deduction based on the clue and feedback, followed immediately by a single line starting with "guess:" containing exactly one valid five{-}letter English word. I will ensure the explanation does not contain conversational filler, self{-}correction, or incomplete thoughts, and I will verify the word length before submitting. Specifically, for the clue "1985 hit by heart," I will identify the song "What About Love" or the band name "Heart" and select a valid 5{-}letter candidate like "ABOUT" or "HEART" (if allowed as a word) or a related term, ensuring the output is clean and strictly follows the `explanation:` and `guess:` tags without extra text.} \\
            }
        }
    }
    & \\ \\

\end{supertabular}
}

\end{document}
